\documentclass[
    oneside, % no mirrored margin for twosided book
    % draft, % no figure, no listings, draw Bold where box overflow (debug)
    12pt % fontsize
]{article}

\usepackage[]{guap}

\usepackage{tikz}
\usetikzlibrary{arrows.meta, positioning}

\usepackage[
    bookmarks
]{hyperref} % Обязательно в конце

\graphicspath{{img/}}

\begin{document}

\newcommand{\shot}[2]{
    \begin{figure}[H]
        \centering
        \fbox{\parbox{0.85\textwidth}{\centering\vspace{4cm}\textit{Скриншот: #1}\vspace{4cm}}}
        \caption{#2}
    \end{figure}
}

\section{Постановка задачи}
Общее задание: меню пользователя должно содержать возможность редактирования любой
записи любого объекта, удаления любого объекта, добавления нового объекта на любую
позицию.

Вариант 7.

\textbf{Задача 1.} Определить класс с именем TRAIN, содержащий следующие поля:
название пункта назначения; номер поезда; время отправления. Определить методы
доступа к этим полям и перегруженные операции извлечения и вставки для объектов
типа Train. Заранее число объектов не известно. Написать программу, выполняющую
следующие действия:
\begin{itemize}
    \item записи должны быть размещены в алфавитном порядке по названиям пунктов
    назначения;
    \item вывод на экран информации о поездах, отправляющихся после введенного с
    клавиатуры времени;
    \item если таких поездов нет, выдать на дисплей соответствующее сообщение.
\end{itemize}

\textbf{Задача 2.} С использованием файловых и строковых потоков написать программу,
которая считывает текст из файла и определяет, сколько в нем слов, состоящих не более
чем из четырех букв.

Ссылка на проект: \url{https://github.com/John-ik/tech_prog_lab2.git}

\section{Формализация}
\textbf{Задача 1.} Класс TRAIN хранит пункт назначения (\texttt{destination}), номер
поезда (\texttt{number}) и время отправления в виде часов и минут (\texttt{hour},
\texttt{minute}). Время вводится в формате ЧЧ:ММ, допустимый диапазон 00:00--23:59.
Для сравнения времени отправления $t$ с введённым временем $T$ используется число минут
от начала суток: $t = 60 \cdot \texttt{hour} + \texttt{minute}$. Поезд попадает в
выборку, если $t > T$.

Список TrainList хранит динамический массив указателей \texttt{Train **} в куче.
Каждый поезд создаётся через \texttt{new} и удаляется списком. При заполнении массив
увеличивается вдвое. Сортировка выполняется вставками по полю \texttt{destination}
по возрастанию.

\textbf{Задача 2.} Текст читается построчно (\texttt{std::getline}). Каждая строка
разбивается на слова по пробельным символам (\texttt{std::istringstream}). Буквами
считаются латинские буквы и русские буквы (А--я, Ё, ё) в кодировке UTF-8. Слово
учитывается, если в нём от 1 до 4 букв. Слова без букв (например, «123») не
учитываются, дефис внутри слова не разделяет его.

\section{Таблица описания классов}
\begin{table}[H]
    \caption{Описание классов}
    \begin{tabular}{|p{3.2cm}|p{6.8cm}|p{4.4cm}|}
        \hline
        Класс & Назначение & Наследование \\
        \hline
        Train & Данные одного поезда: пункт назначения, номер, время отправления;
        методы доступа, ввод и вывод & Нет \\
        \hline
        TrainList & Динамический список поездов: вставка, удаление, замена,
        сортировка; владеет объектами Train & Нет. Содержит Train (композиция) \\
        \hline
        TrainError & Исключение с текстом, поясняющим ошибку ввода или обращения
        к элементу & Наследует \texttt{std::runtime\_error}, тип public \\
        \hline
    \end{tabular}
    \par\smallskip
    Примечание: TrainList удаляет все объекты TRAIN в своём деструкторе, поэтому
    владение объектами не передаётся наружу.
\end{table}

На рисунке показаны связи классов: наследование (TrainError от
\texttt{std::runtime\_error}) и композиция (TrainList содержит Train).

\begin{figure}[H]
    \centering
    \begin{tikzpicture}[node distance=1.6cm, every node/.style={font=\small}]
        \node[draw, minimum width=4.2cm, minimum height=0.8cm] (rt) {\texttt{std::runtime\_error}};
        \node[draw, minimum width=4.2cm, minimum height=0.8cm, below=of rt] (te) {TrainError};
        \node[draw, minimum width=4.2cm, minimum height=0.8cm, below=3cm of te] (tl) {TrainList};
        \node[draw, minimum width=4.2cm, minimum height=0.8cm, right=5cm of tl] (tr) {Train};

        \draw[-{Triangle[open]}, thick] (te) -- (rt);
        \draw[{Diamond[length=10pt]}-, thick] (tl.east) -- (tr.west)
            node[midway, above] {0..*};
    \end{tikzpicture}
    \caption{Иерархия классов: наследование и композиция}
\end{figure}

\section{Исходный код}

\newcommand{\mylisting}[1]{
    \stepcounter{listing}
    \noindent\small
    Листинг \thelisting~--- \texttt{\detokenize{#1}}
    \inputminted[fontsize=\small]{C++}{#1}
}

\mylisting{main1.cpp}
\mylisting{./include/exceptions.hpp}
\mylisting{./include/input.hpp}
\mylisting{./include/log.hpp}
\mylisting{./include/train.hpp}
\mylisting{./include/train_list.hpp}
\mylisting{./src/train.cpp}
\mylisting{./src/train_list.cpp}
\mylisting{main2.cpp}

\section{Результаты работы программы}
Программа задачи 1 запускается из каталога \texttt{build}: \texttt{./task1}.
Команды меню вводятся латинскими буквами: \texttt{p} --- показать список, \texttt{a} ---
добавить в конец, \texttt{i} --- добавить на позицию, \texttt{e} --- редактировать,
\texttt{d} --- удалить, \texttt{s} --- отсортировать по пункту назначения, \texttt{t} ---
поезда, отправляющиеся после введённого времени, \texttt{h} --- помощь, \texttt{q} --- выход.
Помощь выводится по команде \texttt{h} и при неизвестной команде.
Тестовые данные: \texttt{test/task1\_valid.txt} (корректный ввод) и
\texttt{test/task1\_errors.txt} (ошибочный ввод). Файл подаётся на стандартный ввод:
\texttt{./task1 < test/task1\_valid.txt}.

\shot{добавление трёх поездов и вывод списка}{Добавление поездов в конец и вывод списка}
Пользователь трижды вводит команду \texttt{a} и заполняет пункт назначения, номер и время
отправления. Программа выводит сообщения о вызове конструкторов и деструкторов TRAIN, а
команда \texttt{p} показывает поезда в порядке добавления.

\shot{сортировка и поиск по времени}{Сортировка по пункту назначения и поиск поездов после времени}
Команда \texttt{s} упорядочивает список по пункту назначения. Затем команда \texttt{t} и время
12:00 выводит только поезда, отправляющиеся позже 12:00. Если такого поезда нет, выводится
сообщение «Поездов, отправляющихся после ..., нет».

\shot{вставка на позицию, редактирование, удаление}{Вставка на позицию, редактирование и удаление записей}
Пользователь вводит команду \texttt{i}, указывает позицию 2 и поезд; затем \texttt{e} ---
редактирование позиции 1 и \texttt{d} --- удаление позиции 3. Каждая операция сопровождается
выводом списка через \texttt{p}.

\shot{ошибочный ввод}{Обработка ошибок ввода}
Файл \texttt{test/task1\_errors.txt} содержит неизвестную команду \texttt{x}, пустой пункт
назначения, время 25:00, позицию вне диапазона, нецелое значение \texttt{abc} в поле и
неверный формат времени \texttt{1200}. Программа выводит «Ошибка: ...» с пояснением, при
неизвестной команде дополнительно выводит помощь, и продолжает работу.

Задача 2. Программа запрашивает имя файла, выводит число слов из 1--4 букв.
Тестовый файл: \texttt{test/text2.txt}.

\shot{подсчёт слов в тестовом файле}{Подсчёт слов из не более чем четырёх букв}
Для файла \texttt{test/text2.txt} программа выводит «Слов, состоящих не более чем из
четырёх букв: 24». Число совпадает с ручным подсчётом по тексту. При отсутствии файла
выводится «Ошибка: не удалось открыть файл» и программа завершается с кодом 1.

\section{Выводы}
В ходе работы создан класс TRAIN с методами доступа и перегруженными операторами
извлечения и вставки. Список поездов реализован на динамическом массиве указателей без
STL-контейнеров. Ошибки обрабатываются исключением TrainError с текстом. Для задачи 2
использованы файловые и строковые потоки. Конструкторы и деструкторы выводят своё
сообщение.

Ограничения: сортировка выполняется по команде меню, а не при каждой вставке.
Сравнение пунктов назначения побайтовое (UTF-8): учитывается регистр, латиница стоит
перед кириллицей, а «Ё» --- перед «А». Это не совпадает с алфавитом русского языка.

\end{document}
